% Lösungen Einheit 2
\einheitenkopf{Einheit 2 von 4 · Normalform und p-q-Formel}
\erg{11}{a) p-q-Formel \quad b) Wurzelziehen \quad c) Rückwärtsrechnen \quad d) Wurzelziehen \quad e) p-q-Formel (vorher durch $2$ teilen)}
\erg{12}{a) $p = 7$, $q = 10$ \quad b) $p = -3$, $q = 2$ \quad c) $p = 4$, $q = -7$ \quad d) $p = -1$, $q = -12$ \quad e) $p = -11$, $q = 0$}
\erg{13}{a) Zahl vor $x^2$ ist eins \quad b) teilen durch $2$ \quad c) teilen durch $5$ \quad d) teilen durch $-1$ \quad e) teilen durch $0{,}5$}
\erg{14}{a) $-3 \pm 1$; $x_1 = -2$, $x_2 = -4$ \quad b) $-5 \pm 4$; $x_1 = -1$, $x_2 = -9$ \quad c) $-3 \pm 2$; $x_1 = -1$, $x_2 = -5$ \quad d) $-5 \pm 3$; $x_1 = -2$, $x_2 = -8$ \quad e) $3 \pm 1$; $x_1 = 4$, $x_2 = 2$ \quad f) $-1 \pm 5$; $x_1 = 4$, $x_2 = -6$}
\erg{15}{a) $x^2 - 4x - 5 = 0$; $x_1 = 5$, $x_2 = -1$ \quad b) $x^2 + 4x - 12 = 0$; $x_1 = 2$, $x_2 = -6$ \quad c) $x^2 - 4x + 3 = 0$; $x_1 = 3$, $x_2 = 1$ \quad d) $x^2 - 8x + 12 = 0$; $x_1 = 6$, $x_2 = 2$}
\erg{16}{a) $-1{,}5 \pm 3{,}5$; $x_1 = 2$, $x_2 = -5$ \quad b) Wert unter der Wurzel $0$; $x = 5$ \quad c) Wert unter der Wurzel $-3$; keine Lösung \quad d) $6{,}25 - 3 = 3{,}25$; zwei Lösungen \quad e) $x_{1,2} = 2 \pm \sqrt{3}$; $x_1 \approx 3{,}73$, $x_2 \approx 0{,}27$ \quad f) $-1 \pm 6$; $x_1 = 5$, $x_2 = -7$; Probe: $5^2 + 2 \cdot 5 - 35 = 0$ (wA)}
\erg{17}{a) $x^2 - x - 6 = 0$; $x_1 = 3$, $x_2 = -2$ \quad b) $x^2 + 4x - 21 = 0$; $x_1 = 3$, $x_2 = -7$ \quad c) Wurzelziehen: $x^2 = 25$; $x_1 = 5$, $x_2 = -5$ \quad d) Rückwärtsrechnen: $x - 1 = 4$ oder $x - 1 = -4$; $x_1 = 5$, $x_2 = -3$ \quad e) $f(x) = 0$, durch $2$ teilen: $x^2 - 6x + 5 = 0$; $x_1 = 5$, $x_2 = 1$}
\erg{18}{a) $A({-3}|4)$, $B(2|{-1})$ \quad b) $x^2 + x - 6 = 0$; $x_1 = 2$, $x_2 = -3$; Schnittpunkte $(2|{-1})$ und $({-3}|4)$ \quad c) $x^2 - 3x - 4 = 0$; $x_1 = 4$, $x_2 = -1$; $k(4) = 6$, $k(-1) = 1$; Schnittpunkte $(4|6)$ und $({-1}|1)$}
\erg{19}{a) Nicht normiert – erst durch $2$ teilen: $x^2 - 5x + 6 = 0$; $2{,}5 \pm 0{,}5$; $x_1 = 3$, $x_2 = 2$ \quad b) $x^2 + 2x - 3 = 0$; $x_1 = 1$, $x_2 = -3$ \quad c) Vorzeichen von $-\frac{p}{2}$ falsch, richtig ist $+5$: $5 \pm 3$; $x_1 = 8$, $x_2 = 2$ \quad d) $7 \pm 3$; $x_1 = 10$, $x_2 = 4$}
\erg{20}{a) Unter der Wurzel steht $1 - 5 = -4$; aus einer negativen Zahl gibt es keine Wurzel. \quad b) Die Formel gilt nur für die Normalform mit der Zahl eins vor $x^2$; erst durch $4$ teilen: $x^2 + 2x - 1 = 0$, also $p = 2$, $q = -1$. \quad c) Eine Lösung: unter der Wurzel steht $9 - 9 = 0$ (es ist $(x - 3)^2 = 0$, $x = 3$).}
\erg{21}{a) $0{,}01v^2 + 0{,}3v = 40$, durch $0{,}01$ teilen: $v^2 + 30v - 4000 = 0$ \quad b) $v_{1,2} = -15 \pm \sqrt{225 + 4000} = -15 \pm 65$; $v_1 = 50$, $v_2 = -80$ \quad c) Eine Geschwindigkeit ist nicht negativ, also $v = 50\,$km/h; das Auto war $20\,$km/h zu schnell.}
